% !TEX root = ../../main_docker.tex
% Dockerfile

\section{Dockerfile}
\label{sec:dockerfile}

\subsection{Dockerfile: Aufbau und Build-Kontext}
\begin{ztab}{@{}lX@{}}
Format & \R{INSTRUKTION argument} -- Instruktion GROSS, Argumente klein \\
Erste Zeile & muss \R{FROM} sein \\
\R{\textbackslash} & Zeilenumbruch, wird vor der Ausführung entfernt (wie in C) \\
\R{\#} & Kommentar, wird ebenfalls vorher entfernt \\
\end{ztab}

\SW{Build-Kontext} \R{docker build <kontext>} -- das Verzeichnis, das an
den Daemon geht. Docker sucht dort die Datei \R{Dockerfile}, und der
Build-Prozess kann auf \emph{alle} Dateien des Kontexts zugreifen.
\begin{lstlisting}
docker build . -t my_image                  # Standardfall
docker build subdir1 -t my_image            # Kontext woanders
docker build . -f Dockerfile.v1 -t my_image # anderes Dockerfile
docker build subdir1 -f ../subdir2/Dockerfile.v1 -t my_image
\end{lstlisting}

\begin{falle}[„no such file or directory: Dockerfile“]
Meldet der Build \R{transferring dockerfile: 2B}, liegt im Kontext keine
Datei namens \R{Dockerfile} -- grosses D, Rest klein, \emph{keine} Endung
(Linux ist case-sensitiv, \R{Dockerfile.txt} ist der Windows-Klassiker).
Prüfen mit \R{ls -lA}.\\
\R{-f} bestimmt nur die \emph{Bauanleitung}; der Kontext bleibt das
Verzeichnis am Schluss. Alle \R{COPY}-Pfade sind relativ zum
\emph{Kontext}, nicht zum Ort des Dockerfiles.
\end{falle}

\begin{formel}[Was beim Bauen passiert]
Jeder Schritt erzeugt einen temporären Container-Layer und speichert ihn
dann als Image ab -- technisch wie \R{docker commit}. Das zuletzt so
erzeugte Image ist das Ergebnis. Wenn davon die Rede ist, „etwas im Image
zu tun“, ist immer der Container gemeint, der das Image ergibt.
\end{formel}

\subsection{FROM, COPY, RUN}
\SW{\R{FROM <image>}} legt das \T{Parent-Image} fest. Parent-Images haben
selbst wieder Parents -- daher die Schichten:
\R{kaohslu/epp-demo} $\to$ .NET-7.0-Runtime $\to$ Debian.
Offizielle Images gibt es für Linux-Distributionen (\T{Busybox, Alpine,
Debian, Ubuntu} -- von klein nach gross), Entwicklungsumgebungen
(.NET, Python, GCC, Node.js) und OSS (MySQL, InfluxDB, Grafana, WordPress,
Node-Red, Nginx) -- letztere enthalten selbst eine Linux-Distribution.

\SW{\R{COPY <src> <dest>}} Quellen liegen im Build-Kontext des Hosts,
absolut oder relativ dazu; ganze Verzeichnisse gehen auch. Mehrere Quellen
nur, wenn \R{<dest>} ein Verzeichnis ist. Das Ziel ist absolut oder
relativ zum \R{WORKDIR}.
\begin{ztab}{@{}lX@{}}
\R{COPY ["<src>","<dest>"]} & Klammerform, falls Pfade Leerzeichen haben \\
\R{--chown=<user>:<group>} & Eigentümer im Image setzen \\
\R{--chmod=<perms>} & Berechtigungen setzen \\
\R{ADD} & wie \R{COPY}, aber mit „Magie“ (packt Archive automatisch aus)
  \ra nur verwenden, wenn man genau das braucht \\
\end{ztab}

\SW{\R{RUN <cmd>}} führt Befehle im entstehenden Image aus: Pakete
installieren, Compiler oder Archiv-Werkzeuge aufrufen. Die
Dateisystemänderungen stehen dem nächsten Schritt zur Verfügung.
\begin{ztab}{@{}lX@{}}
\R{RUN echo "Hallo" > greeting.txt} & wirkt -- ändert das Dateisystem \\
\R{RUN echo "Hallo"} & \textbf{wirkungslos} -- nur Konsolenausgabe, keine
  Dateisystemänderung, bleibt im Image ohne Spur \\
\end{ztab}
In Dockerfiles nimmt man \T{\R{apt-get}} statt \R{apt} -- \R{apt-get} ist
für Skripte gedacht, \R{apt} für Nutzerinteraktion.

\begin{falle}[Shell-Form gegen Exec-Form]
\R{RUN echo "x" > /tmp/f} läuft über \R{/bin/sh -c}, eine Shell
interpretiert also das \R{>}. \R{RUN ["echo","x",">","/tmp/f"]} startet
\R{echo} \emph{ohne} Shell -- \R{>} wird zum blossen Argument und die
Datei entsteht nie. Bei \R{CMD} und \R{ENTRYPOINT} ist es umgekehrt: dort
nimmt man die Klammerform.
\end{falle}

\subsection{Schichten im Image}
\label{rez:layer}
\begin{formel}[Die Regel]
Jede Zeile im Dockerfile, \emph{die das Dateisystem ändert}, erzeugt einen
eigenen Image-Layer. \R{CMD}, \R{ENTRYPOINT}, \R{ENV} und \R{LABEL}
schreiben nur Metadaten \ra 0\,B, kein Layer. \R{WORKDIR} legt ein
Verzeichnis an \ra doch ein Layer.
\end{formel}
\SW{Anzeigen} \R{docker inspect <image>}, Abschnitt \R{RootFS.Layers}
(Anzahl und Hashes) $\cdot$ \R{docker history <image>} (die Instruktion
hinter jeder Schicht und ihre Grösse).
Debian hat nur einen Layer, Ubuntu zwei.

\begin{bsp}[\R{nanodeb}: 3 Layer werden 2]
\begin{lstlisting}
FROM debian                    FROM debian
RUN apt-get update             RUN apt-get update && \
RUN apt-get install -y nano        apt-get install -y nano && \
                                   apt-get clean && \
                                   rm -rf /var/lib/apt/lists/*
=> 3 Layer, 140 MB             => 2 Layer, 122 MB
\end{lstlisting}
Das \R{update} hinterlässt unnötige Daten in einer eigenen Schicht. Im
selben \R{RUN} aufgeräumt, verschwinden sie -- in einem \emph{späteren}
\R{RUN} gelöscht, blieben sie in der Schicht darunter liegen.
\end{bsp}

\SW{Best Practice} Zusammengehörige Kommandos in \emph{ein} \R{RUN} mit
\R{\&\&} $\cdot$ unnötige Dateien im selben Schritt löschen
(\R{apt-get clean}, \R{rm -rf /var/lib/apt/lists/*}) $\cdot$ sinnvolle
Reihenfolge der Instruktionen, damit der \T{Build-Cache} möglichst lange
greift. Eine geänderte Schicht macht alle darunterliegenden ungültig.

\subsection{Startkommando: \R{CMD} und \R{ENTRYPOINT}}
\label{form:entrypoint}
\R{docker run <image> [command] [args...]} -- was gestartet wird, steht im
Image unter \R{Config.Cmd} und \R{Config.Entrypoint}
(\R{docker inspect}). Nützliches Alias:
\begin{lstlisting}
alias dentry="docker inspect -f \
  'ENTRYPOINT:{{.Config.Entrypoint}}; CMD:{{.Config.Cmd}}'"
\end{lstlisting}
So hat \R{ubuntu} ein leeres ENTRYPOINT und \R{CMD ["/bin/bash"]} --
deshalb startet dort ohne weitere Angabe eine \R{bash}.
\R{kaohslu/epp-demo} hat \R{ENTRYPOINT ["dotnet","ASYD\_Demo.dll"]} und
leeres CMD. Prüfen lässt sich das Ergebnis mit \R{docker top <cnt>}.

\begin{formel}[Der entscheidende Unterschied]
\R{CMD} gesetzt: \R{[command]} wird \emph{anstelle} von CMD ausgeführt.\\
\R{ENTRYPOINT} gesetzt: \R{[command] [args...]} werden \emph{angehängt}.\\
Beide gesetzt: ENTRYPOINT ist das Kommando, CMD liefert die
Standard-Argumente dafür.
\begin{tabularx}{\linewidth}{@{}llll@{}}
ENTRYPOINT & CMD & \R{[command]} & ausgeführt wird \\
\R{["script.sh"]} & -- & -- & \R{script.sh} \\
\R{["script.sh"]} & -- & \R{/bin/bash} & \R{script.sh /bin/bash} \\
\R{["script.sh"]} & \R{["mysqld"]} & -- & \R{script.sh mysqld} \\
\R{["script.sh"]} & \R{["mysqld"]} & \R{/bin/bash} & \R{script.sh /bin/bash} \\
-- & \R{["/bin/sh"]} & -- & \R{/bin/sh} \\
-- & \R{["/bin/sh"]} & \R{/bin/bash} & \R{/bin/bash} \\
\end{tabularx}
\end{formel}

\SW{\R{WORKDIR <pfad>}} setzt das Arbeitsverzeichnis für alle folgenden
\R{RUN}, \R{CMD}, \R{ENTRYPOINT}, \R{COPY} und \R{ADD}. Es wird angelegt,
auch wenn es gar nicht benutzt wird. Entspricht dem, was \R{pwd} im
Container ausgibt.

\begin{bsp}[Die \R{nanodeb}-Reihe]
\begin{lstlisting}
v3:  CMD ["nano","/greeting.txt"]
     run ohne Arg  -> nano /greeting.txt
     run mit bash  -> ersetzt CMD, startet bash
v4:  ENTRYPOINT ["nano"]  CMD ["/greeting.txt"]
     run ohne Arg  -> nano /greeting.txt
     run mit main.c-> nano main.c
v5:  + WORKDIR /app
     run mit main.c-> nano /app/main.c
\end{lstlisting}
\end{bsp}
